\textbf{Phần 1. Trắc nghiệm lựa chọn}
\setcounterpageref{ex}{0}

\begin{ex} %[1D6V23]
Trong bảng các số từ 1 đến 100, số đứng ngay sau số 99 là số nào?
\choice
{$90$}
{$110$}
{\True $100$}
{$99$}
\loigiai{
Số đứng liền sau số 99 là số $100$.
\textbf{Đáp số:} $100$.\par
\textbf{Đáp án của câu hỏi là C}
}
\end{ex}

\begin{ex} %[1D6V23]
Trong bảng các số từ 1 đến 100, có bao nhiêu số có hai chữ số?
\choice
{$80$}
{\True $90$}
{$100$}
{$89$}
\loigiai{
Các số có hai chữ số từ 10 đến 99 có: $99 - 10 + 1 = 90$ (số).
\textbf{Đáp số:} $90$.\par
\textbf{Đáp án của câu hỏi là B}
}
\end{ex}

\begin{ex} %[1D6V23]
Trong bảng các số từ 1 đến 100, có bao nhiêu số có một chữ số?
\choice
{$8$}
{$10$}
{$7$}
{\True $9$}
\loigiai{
Từ 1 đến 100 có các số một chữ số là: 1, 2, 3, 4, 5, 6, 7, 8, 9 (tất cả 9 số).
\textbf{Đáp số:} $9$.\par
\textbf{Đáp án của câu hỏi là D}
}
\end{ex}

\begin{ex} %[1D6V23]
Số bé nhất có ba chữ số trong bảng là số nào?
\choice
{\True $100$}
{$90$}
{$110$}
{$99$}
\loigiai{
Số 100 là số duy nhất có ba chữ số trong bảng các số từ 1 đến 100.
\textbf{Đáp số:} $100$.\par
\textbf{Đáp án của câu hỏi là A}
}
\end{ex}

\begin{ex} %[1D6V23]
Trong cùng một cột của bảng số từ 1 đến 100, hai số ở hai hàng liền nhau hơn kém nhau bao nhiêu đơn vị?
\choice
{$0$}
{$20$}
{\True $10$}
{$9$}
\loigiai{
Trong bảng số, mỗi hàng có 10 số, do đó hai ô ở hai hàng liền nhau trên cùng một cột hơn kém nhau đúng $10$ đơn vị.
\textbf{Đáp số:} $10$.\par
\textbf{Đáp án của câu hỏi là C}
}
\end{ex}

\begin{ex} %[1D6V23]
Số ở ô ngay phía dưới số 45 trong bảng số là số nào?
\choice
{$54$}
{\True $55$}
{$56$}
{$53$}
\loigiai{
Số ở ô ngay dưới cùng cột cách số ở trên 10 đơn vị: $45 + 10 = 55$.
\textbf{Đáp số:} $55$.\par
\textbf{Đáp án của câu hỏi là B}
}
\end{ex}

\begin{ex} %[1D6V23]
Số ở ô ngay phía trên số 68 trong bảng số là số nào?
\choice
{$57$}
{$59$}
{$56$}
{\True $58$}
\loigiai{
Số ở ô ngay trên cùng cột kém số ở dưới 10 đơn vị: $68 - 10 = 58$.
\textbf{Đáp số:} $58$.\par
\textbf{Đáp án của câu hỏi là D}
}
\end{ex}

\begin{ex} %[1D6V23]
Số ở ô ngay bên phải của số 37 trong bảng số là số nào?
\choice
{\True $38$}
{$37$}
{$39$}
{$36$}
\loigiai{
Số ở ngay bên phải là số liền sau: $37 + 1 = 38$.
\textbf{Đáp số:} $38$.\par
\textbf{Đáp án của câu hỏi là A}
}
\end{ex}

\begin{ex} %[1D6V23]
Số ở ô ngay bên trái của số 80 trong bảng số là số nào?
\choice
{$78$}
{$80$}
{\True $79$}
{$77$}
\loigiai{
Số ở ngay bên trái là số liền trước: $80 - 1 = 79$.
\textbf{Đáp số:} $79$.\par
\textbf{Đáp án của câu hỏi là C}
}
\end{ex}

\begin{ex} %[1D6V23]
Hàng thứ 5 trong bảng các số từ 1 đến 100 gồm các số từ 41 đến bao nhiêu?
\choice
{$40$}
{\True $50$}
{$60$}
{$49$}
\loigiai{
Hàng thứ 5 gồm 10 số: $41, 42, 43, 44, 45, 46, 47, 48, 49, 50$.
\textbf{Đáp số:} $50$.\par
\textbf{Đáp án của câu hỏi là B}
}
\end{ex}

\begin{ex} %[1D6V23]
Cột cuối cùng bên phải của bảng số gồm toàn các số tròn chục. Số cuối cùng ở góc dưới bên phải là số nào?
\choice
{$90$}
{$110$}
{$99$}
{\True $100$}
\loigiai{
Số ở góc dưới cùng bên phải của bảng số là số $100$.
\textbf{Đáp số:} $100$.\par
\textbf{Đáp án của câu hỏi là D}
}
\end{ex}

\begin{ex} %[1D6V23]
Có bao nhiêu số có chữ số tận cùng là 7 trong bảng các số từ 1 đến 100?
\choice
{\True $10$}
{$0$}
{$20$}
{$9$}
\loigiai{
Các số đó là: $7, 17, 27, 37, 47, 57, 67, 77, 87, 97$. Có tất cả $10$ số.
\textbf{Đáp số:} $10$.\par
\textbf{Đáp án của câu hỏi là A}
}
\end{ex}

\textbf{Phần 2. Trắc nghiệm đúng sai}
\setcounterpageref{ex}{0}

\begin{ex} %[1D6V23]
Xét tính đúng hoặc sai của các khẳng định toán học sau đây:
\choiceTF
{\True Kết quả của bài toán «Có bao nhiêu số có chữ số hàng chục là 5 trong bảng các số từ 1 đến 100» là $10$.}
{Kết quả của bài toán «Tìm số đứng liền trước của số 100» là $98$.}
{Kết quả của bài toán «Tìm số tự nhiên $x$ biết $x$ đứng giữa số 88 và số 90 trong bảng số» là $88$.}
{\True Đếm cách 5 đơn vị: $5, 10, 15, 20, \dots, 35$. Số thích hợp điền vào chỗ chấm là số $25$.}
\loigiai{
\begin{itemize}
    \item Ý a) Đúng: Các số đó là hàng từ 50 đến 59: $50, 51, \dots, 59$. Có tất cả $10$ số.
\textbf{Đáp số:} $10$.
    \item Ý b) Sai: Kết quả đúng là $99$, không phải $98$. Số liền trước của 100 là: $100 - 1 = 99$.
\textbf{Đáp số:} $99$.
    \item Ý c) Sai: Kết quả đúng là $89$, không phải $88$. Số đứng giữa 88 và 90 là số $89$.
\textbf{Đáp số:} $89$.
    \item Ý d) Đúng: Đếm thêm 5 đơn vị: $20 + 5 = 25$.
\textbf{Đáp số:} $25$.
\end{itemize}
\begin{itemize}
    \item a) \textbf{Đúng}.
    \item b) \textbf{Sai}.
    \item c) \textbf{Sai}.
    \item d) \textbf{Đúng}.
\end{itemize}
}
\end{ex}

\begin{ex} %[1D6V23]
Xét tính đúng hoặc sai của các bài toán và phép tính sau đây:
\choiceTF
{Đếm cách 10 đơn vị: $10, 20, 30, \dots, 60$. Số thích hợp điền vào chỗ chấm là số $41$.}
{\True Kết quả của bài toán «Trong bảng số, từ số 21 đến số 30 có tất cả bao nhiêu số» là $10$.}
{\True Trong bảng số, từ số 1 đến số 50 có $5$ số tròn chục.}
{Kết quả của bài toán «Tìm số tự nhiên $x$ biết $x$ đứng giữa số 88 và số 90 trong bảng số» là $88$.}
\loigiai{
\begin{itemize}
    \item Ý a) Sai: Kết quả đúng là $40$, không phải $41$. Đếm thêm 10 đơn vị: $30 + 10 = 40$.
\textbf{Đáp số:} $40$.
    \item Ý b) Đúng: Số lượng số từ 21 đến 30 là: $30 - 21 + 1 = 10$ (số).
\textbf{Đáp số:} $10$.
    \item Ý c) Đúng: Các số tròn chục là: $10, 20, 30, 40, 50$. Có tất cả $5$ số.
\textbf{Đáp số:} $5$.
    \item Ý d) Sai: Kết quả đúng là $89$, không phải $88$. Số đứng giữa 88 và 90 là số $89$.
\textbf{Đáp số:} $89$.
\end{itemize}
\begin{itemize}
    \item a) \textbf{Sai}.
    \item b) \textbf{Đúng}.
    \item c) \textbf{Đúng}.
    \item d) \textbf{Sai}.
\end{itemize}
}
\end{ex}

\textbf{Phần 3. Trả lời ngắn}
\setcounterpageref{ex}{0}

\begin{ex} %[1D6V23]
Tìm một số biết rằng số đó ở ngay bên dưới số 72 và ở ngay bên trái số 83.
\shortans{82}
\loigiai{
\textbf{Đáp số:} $82$\par
Số ở dưới số 72 là: $72 + 10 = 82$. Số ở bên trái số 83 là: $83 - 1 = 82$. Vậy số đó là $82$.
\textbf{Đáp số:} $82$.
}
\end{ex}

\begin{ex} %[1D6V23]
Một hàng có 10 ô số. Bạn Hùng đã điền được 6 ô số. Hỏi Hùng còn phải điền bao nhiêu ô số nữa để hoàn thành hàng đó?
\shortans{4}
\loigiai{
\textbf{Đáp số:} $4$\par
Số ô số Hùng còn phải điền là: $10 - 6 = 4$ (ô số).
\textbf{Đáp số:} $4$.
}
\end{ex}

\begin{ex} %[1D6V23]
Một bảng số bị rách một góc làm mất các số từ 95 đến 100. Hỏi có bao nhiêu số bị mất?
\shortans{6}
\loigiai{
\textbf{Đáp số:} $6$\par
Số các số bị mất là: $100 - 95 + 1 = 6$ (số).
\textbf{Đáp số:} $6$.
}
\end{ex}

\begin{ex} %[1D6V23]
Số lớn nhất có hai chữ số khác nhau trong bảng số là số nào?
\shortans{98}
\loigiai{
\textbf{Đáp số:} $98$\par
Số lớn nhất có 2 chữ số khác nhau là $98$.
\textbf{Đáp số:} $98$.
}
\end{ex}

\begin{ex} %[1D6V23]
Tổng của chữ số hàng chục và hàng đơn vị của số 96 bằng bao nhiêu?
\shortans{15}
\loigiai{
\textbf{Đáp số:} $15$\par
Tổng hai chữ số là: $9 + 6 = 15$.
\textbf{Đáp số:} $15$.
}
\end{ex}

\begin{ex} %[1D6V23]
Hiệu giữa chữ số hàng chục và hàng đơn vị của số 73 bằng bao nhiêu?
\shortans{4}
\loigiai{
\textbf{Đáp số:} $4$\par
Hiệu hai chữ số là: $7 - 3 = 4$.
\textbf{Đáp số:} $4$.
}
\end{ex}
